\documentclass[12 pt]{article}

\usepackage{stoversymb}
\usepackage[margin=0.75in, top=0.875in, bottom=0.875in]{geometry}
\usepackage{amsmath,amsfonts,amssymb,url,multicol,graphicx,tikz,soul,changepage}
%===makes urls render well===
\usepackage{lmodern}
\usepackage[T1]{fontenc}
%============================
\usepackage{wasysym} % smileys
\usepackage[inline]{enumitem}
\usepackage[many]{tcolorbox}
\tcbuselibrary{breakable}

%\everymath{\displaystyle}

\makeatletter
\renewcommand*\env@matrix[1][\arraystretch]{%
	\edef\arraystretch{#1}%
	\hskip -\arraycolsep
	\let\@ifnextchar\new@ifnextchar
	\array{*\c@MaxMatrixCols c}}
\makeatother

%\setenumerate{itemsep=0.25in}
%\setlist[enumerate,1]{label=\arabic*., itemsep=0.5in}
%\setlist[enumerate,2]{label={(\alph*)}, itemsep=0.8in}
%\setlist[enumerate,2]{leftmargin=0.5in}
%\setlist[enumerate,3]{label={\roman*.}}
\setlist[itemize,1]{label=$\circ$,topsep=0mm,itemsep=0mm}

\newcommand{\truefalse}[1]{#1\hfill\rule[-1mm]{220pt}{0.75pt}}
\newcommand{\hint}[1]{\vspace{3mm}\textbf{Hint}: #1}
\newcommand{\note}[1]{\textbf{Note}: #1}
\newcommand{\infsum}[3]{\sum_{{#1}={#2}}^\infty {#3}}
\newcommand{\axes}[1]{\begin{center}\includegraphics[scale=#1]{3DAxes}\end{center}}
\newcommand{\comps}[1]{\langle #1_1,#1_2,#1_3\rangle}
\newcommand{\compslong}[3]{\langle #1, #2, #3\rangle}
\newcommand{\ijk}[2]{#1\vect{#2}}
%\newcommand{\pmat}[1]{\begin{pmatrix}[1.25]#1\end{pmatrix}}
%\newcommand{\bmat}[1]{\begin{bmatrix}[1.25]#1\end{bmatrix}}
\newcounter{Exx}
\newcommand{\exbox}[1]{\stepcounter{Exx}\begin{tcolorbox}[breakable,arc=2pt,boxrule=1pt,right=6mm,top=3mm,bottom=6mm]\noindent\textit{\large\bfseries Example \theExx:\\[3mm]}~#1\end{tcolorbox}}
\newcommand{\exit}[1]{\vspace{4.5mm}\noindent\textit{Example #1:}}
\newcommand{\exbf}[1]{\vspace{1.5mm}\noindent\textbf{Example #1:}}
\newcommand{\LHS}{\text{LHS}}
\newcommand{\RHS}{\text{RHS}}
\newcommand{\ans}{\noindent\underline{Ans:}~}

\newenvironment{mypmatrix}[1]{\renewcommand{\arraystretch}{#1}\begin{pmatrix}}{\end{pmatrix}}
\newenvironment{mybmatrix}[1]{\renewcommand{\arraystretch}{#1}\begin{bmatrix}}{\end{bmatrix}}
\newcommand{\justify}{\ul{Justify your claim.}}

\newcommand{\pmat}[1]{\begin{mypmatrix}{1.25}#1\end{mypmatrix}}
\newcommand{\bmat}[1]{\begin{mybmatrix}{1.25}#1\end{mybmatrix}}

\makeatletter
\renewcommand*\env@matrix[1][*\c@MaxMatrixCols c]{%
	\hskip -\arraycolsep
	\let\@ifnextchar\new@ifnextchar
	\array{#1}}
\makeatother

\newcommand{\pmatgrid}[2]{\renewcommand{\arraystretch}{1.25}\begin{pmatrix}[#1] #2\end{pmatrix}}
\newcommand{\bmatgrid}[2]{\renewcommand{\arraystretch}{1.25}\begin{bmatrix}[#1] #2\end{bmatrix}}
\newcommand{\warnbox}[1]{\begin{adjustwidth}{0.375in}{0.375in}\begin{tcolorbox}[arc=0pt, colback=red!5!white, colframe=red, boxrule=0.5pt, parbox=false]{\parskip=3mm \parindent=0mm \red{\textbf{\large\itshape Warning!}\par #1}}\end{tcolorbox}\end{adjustwidth}}
\newcommand{\greentheorembox}[2]{\begin{adjustwidth}{0.375in}{0.375in}\begin{tcolorbox}[arc=0pt, colback=darkgreen!10!white, colframe=darkgreen, boxrule=0.5pt, parbox=false]{\parskip=3mm \parindent=0mm \green{\textbf{\large\itshape #1:}\par #2}}\end{tcolorbox}\end{adjustwidth}}
\newcounter{FactBox}
\newcommand{\factbox}[1]{\stepcounter{FactBox}\begin{adjustwidth}{0.375in}{0.375in}\begin{tcolorbox}[arc=0pt, colback=darkgreen!10!white, colframe=darkgreen, boxrule=0.5pt, parbox=false]{\parskip=3mm \parindent=0mm \green{\textbf{\large\itshape Fact \theFactBox:}\par #1}}\end{tcolorbox}\end{adjustwidth}}

\newcommand{\arrow}[1]{\xrightarrow{\hspace*{#1}}}

\setlength\parskip{0.75mm}	

\newcounter{ThmBox}
\newcommand{\thmbox}[1]{\stepcounter{ThmBox}\vspace{6mm}\noindent\fbox{\parbox{\linewidth}{~~(\theThmBox)~\hfill #1 \hfill\hfill}}}
\newcommand{\why}[1]{\begin{adjustwidth}{0.5in}{0.5in}{\par\vspace{3mm}\noindent\green{\textbf{Why should this be true? }#1}}\qed\end{adjustwidth}\vspace{6mm}}

\graphicspath{ {./../img/} }
\DeclareGraphicsExtensions{.pdf}

\begin{document}
%\begin{flushright}Name: \line(1,0){200}\end{flushright}
\begin{center}
\Large{\textbf{The Inverse Matrix Theorem I}}
\end{center}
Recall that the \ul{inverse} of an $n\times n$ matrix $\sansA$ is an $n\times n$ matrix $\sansA^{-1}$ for which 
\[
	\sansA\sansA^{-1}=\sansI_n=\sansA^{-1}\sansA,
\] 
where $\sansI_n$ is the $n\times n$ identity matrix. Not all matrices have inverses, and those that do are called \ul{invertible} or \ul{nonsingular}.

In general, a matrix being invertible/nonsingular tells you a \textit{tremendous amount} about the matrix + its corresponding linear system(s) / linear transformation(s). To highlight this, your textbook regularly adds to + revisits a monolithic colossus it calls \textbf{The Invertible Matrix Theorem}. 

This so-called ``theorem'' is \ul{really} just a collection of statements/observations which mean the same thing as (and hence are logically equivalent to) \textit{$\sansA$ has an inverse}. However, because many of the statements lumped into this ``theorem'' are important---and indeed, many are related to / duplicates of statements we've already visited before---I want to make sure you have them explicitly given and explained to you. Hence, this handout!

As we learned in class, $\sansA$ has an inverse if and only if $\det(\sansA)\neq 0$; as it turns out, this determinant condition is the very first item in our ``theorem.'' \red{\textbf{Throughout, assume that $\sansA=\pmatgrid{c|c|c}{\vect{v}_1 & \cdots & \vect{v}_n}$ is an $n\times n$ matrix, i.e. that $T(\vect{x})=\sansA\vect{x}$ is a linear transformation $\Reals^n\to\Reals^n$!}}


%==========================================================================
\thmbox{$\sansA$ is invertible if and only if $\det(\sansA)\neq 0$.}
\why{When you think of a matrix being nonsingular, you should think of a linear transformation which neither ``squishes'' nor ``blows up'' information. 

This is possible if and only if such a transformation doesn't transform parallelograms with positive area into those with zero area, and because 
\[
	\text{(area after transformation)} = |\det(\sansA)|\cdot\text{(area before transformation)}\neq 0,
\]
it follows that a matrix is nonsingular if and only if its determinant is not zero.}


%==========================================================================
\vspace{3mm}\noindent A number of other conditions are related to having nonzero determinant. For example:

\thmbox{$\sansA$ is invertible if and only if the columns of $\sansA$ form a linearly independent set.}
\why{A matrix which has linearly dependent rows or columns will result in a cofactor expansion which equals zero (\textbf{check this!}).
	
Thus, $\det(\sansA)\neq 0$ if and only if the columns of $\sansA$ form a linearly independent set.}


%==========================================================================
\thmbox{$\sansA$ is invertible if and only if $\sansA\vect{x}=\vect{0}$ has only the trivial solution.}
\why{As we've seen before, $\sansA\vect{x}=\vect{0}$ has nontrivial solutions if and only if the columns of $\sansA$ are linearly \ul{dependent}. Now, refer to (2).}


\newpage


%==========================================================================
\thmbox{$\sansA$ is invertible if and only if the linear transformation $T(\vect{x})=\sansA\vect{x}$ is one-to-one.}
\why{$T(\vect{x})=\sansA\vect{x}$ is one-to-one if and only if $\sansA\vect{x}=\vect{0}$ has only the trivial solution. Now, refer to (3).}


%==========================================================================
Now, we know that general linear transformations may be one-to-one and not onto or vice versa. However, because $\sansA$ is $n \times n$ (i.e. $T:\Reals^n\to\Reals^n$), this is \textit{not} the case and $T$ being one-to-one is possible if and only if $T$ is also onto. 

If you want, you can accept this as fact; however, it also follows from $\sansA$ being invertible:

\thmbox{$\sansA$ is invertible if and only if the linear transformation $T(\vect{x})=\sansA\vect{x}$ is onto.}
\why{$T(\vect{x})=\sansA\vect{x}$ is onto if and only if every vector $\vect{b}$ in (the codomain) $\Reals^n$ is hit by some vector $\vect{x}$ in (the domain) $\Reals^n$. 
	
Because $\sansA=\pmatgrid{c|c|c}{\vect{v}_1 & \cdots & \vect{v}_n}$, this means that $T$ is onto if and only if the collection $\{\vect{v}_1,\ldots,\vect{v}_n\}$ span $\Reals^n$, but a set of $n$ vectors spans $\Reals^n$ if and only if $\{\vect{v}_1,\ldots,\vect{v}_n\}$ are linearly independent. Now, refer to (2).}


%==========================================================================
As a result of (5), there are a number of equivalent conditions which stem from various rewritings of onto.

\thmbox{$\sansA$ is invertible if and only if the equation $\sansA\vect{x}=\vect{b}$ has $\geq 1$ solution for each $\vect{b}\in\Reals^n$.}
\why{The equation $\sansA\vect{x}=\vect{b}$ has $\geq 1$ solution for each $\vect{b}\in\Reals^n$ if and only if $T(\vect{x})=\sansA\vect{x}$ is onto. Now, apply (5).}


%==========================================================================
\thmbox{$\sansA$ is invertible if and only if the columns of $\sansA$ span $\Reals^n$.}
\why{The columns of $\sansA=\pmatgrid{c|c|c}{\vect{v}_1 & \cdots & \vect{v}_n}$ span $\Reals^n$ if and only if $T(\vect{x})=\sansA\vect{x}$ is onto. Now, apply (5).}


%==========================================================================
\vspace{3mm}The next ingredient of our ``theorem'' comes from the process of finding the inverse of a matrix (when it exists).

\thmbox{$\sansA$ is invertible if and only if $\sansA$ is row equivalent to the $n\times n$ identity matrix $\sansI_n$.}
\why{Recall that---when it exists---you can find $\sansA^{-1}$ by forming the augmented matrix $\pmatgrid{c|c}{\sansA & \sansI_n}$ and putting the ``$\sansA$ part'' into RREF. The result is guaranteed to be the augmented matrix 
\[
	\pmatgrid{c|c}{\sansI_n & \sansA^{-1}},
\]
and so $\sansA$ can be transformed into $\sansI_n$ by elementary row operations.}


%==========================================================================
\vspace{3mm}\noindent Another equivalence involves the relationship between $\sansA$ and its transpose $\sansA^\sansT$.

Recall that the \ul{transpose} of an $m\times n$ matrix $\sansB$ is the $n\times m$ matrix $\sansB^\sansT$ whose first row is the first column of $\sansB$, whose second row is the second column of $\sansB$, etc. 

As was shown on a previous handout, $\det(\sansA)=\det(\sansA^\sansT)$ for all square matrices $\sansA$ and \textit{this} is the key to the point at hand?

\thmbox{$\sansA$ is invertible if and only if $\sansA^\sansT$ invertible.}
\why{$\sansA$ is invertible if and only if $\det(\sansA)\neq 0$ (see (1)) and $\det(\sansA)=\det(\sansA^\sansT)$. Hence, $\sansA$ is invertible if and only if $\det(\sansA^\sansT)\neq 0$ if and only if $\sansA^\sansT$ is invertible.}


%==========================================================================
\vspace{3mm}\noindent The above list of conditions is just a small fraction of the number of equivalent ways one can say ``$\sansA$ is invertible,'' and we'll most definitely be revisiting this handout $\geq 1$ time throughout the semester to make updates!

In the meantime, take some time to read through the above and digest everything thoroughly. To help, use the above to work through the following practice problems!
\end{document}